%----------------------------------------------------------------------------------------
% Tailored CV — Wissenschaftliche*r Mitarbeiter*in, Biochemie und Oberflächenfunktionalisierung
% Fraunhofer HHI, Photonic Biosensors, Berlin
% Ref. 84447
%----------------------------------------------------------------------------------------

\documentclass[10pt,a4paper,sans]{moderncv}

\usepackage[utf8]{inputenc}
\usepackage[T1]{fontenc}

\moderncvstyle{classic}
\moderncvcolor{blue}

\usepackage[scale=0.78]{geometry}
\usepackage{ragged2e}

%----------------------------------------------------------------------------------------

\firstname{Kundai Farai}
\familyname{Sachikonye}

\address{Auf der Lichtung 19}{82178 Puchheim, Germany}
\mobile{(+49) 1626386344}
\email{kundai.sachikonye@bitspark.com}
\homepage{github.com/fullscreen-triangle}{github.com/fullscreen-triangle}

%----------------------------------------------------------------------------------------

\begin{document}

\makecvtitle

%----------------------------------------------------------------------------------------
%	EDUCATION
%----------------------------------------------------------------------------------------

\section{Education}

\cventry{2019--2021}{PhD candidate, Computational Lipidomics}{Technische Universität München / Bayerisches Forschungsinstitut}{}{}{
  Mass spectrometry-based lipidomics, high-dimensional spectral data analysis,
  multi-omics statistical modelling. Departed to pursue independent research.
}
\cventry{2018--2019}{MSc Computational Biology}{Universität Tübingen}{}{}{}
\cventry{2014--2017}{MSc Pharmaceutical Biotechnology}{HAW Hamburg}{}{}{
  Protein biochemistry, bioconjugation, immunoassay development,
  biopharmaceutical process design, analytical chemistry.
}
\cventry{2011--2014}{BSc Biochemistry and Cell Biology}{Jacobs University Bremen}{}{}{
  Protein structure and function, cell biology, analytical biochemistry,
  organic chemistry.
}

%----------------------------------------------------------------------------------------
%	EXPERIENCE
%----------------------------------------------------------------------------------------

\section{Experience}

\cventry{2021--present}{Independent AI \& Computational Biology Developer}{Bitspark GmbH}{}{}{
  Optical physics and molecular spectroscopy from first principles; biosensor
  signal analysis; molecular recognition modelling; binding affinity prediction;
  large-scale spectral data pipelines.
}

\cventry{2019--2021}{Computational Lipidomics}{TUM / Bayerisches Forschungsinstitut}{Munich}{}{
  MS-based lipidomics at HPC scale: peak detection, ion annotation (LIPID MAPS,
  LipidBlast, HMDB), multi-omics integration, federated processing.
}

\cventry{2017--2018}{Glycomics and Proteomics}{Universitätsklinikum Hamburg-Eppendorf}{}{}{
  Automated LC-MS/MS and MALDI-TOF pipeline for proteomics and glycomics:
  protein immobilisation on analytical surfaces, sample preparation protocol
  standardisation, reproducibility validation, SOP development.
}

\cventry{2013}{Cancer Biomarker Discovery}{University Medical Center Groningen}{}{}{
  UPLC-MS/MS shotgun proteomics for cancer biomarker discovery: sample digestion
  and preparation protocols, statistical validation against reference standards.
}

%----------------------------------------------------------------------------------------
%	RELEVANT COMPUTATIONAL TOOLS
%----------------------------------------------------------------------------------------

\section{Computational Tools for Biosensor Applications}

\cvitem{lagrangian}{
  \textbf{First-Principles Optical Physics.}
  Resonant Stack theorem proves Beer-Lambert absorption, Kirchhoff circuit law,
  and thermal emission as formally identical operations ($T=1/e=0.3679$ in all
  three representations). Directly applicable to evanescent field coupling and
  waveguide absorption in PIC biosensors; enables first-principles prediction
  of sensing layer refractive index contributions from molecular structure.
  GPU shader telescope: 10/10 benchmarks, $6.9\times10^{-5}$ median error.
  Live: \href{https://space-observatory-xi.vercel.app}{space-observatory-xi.vercel.app}.
  \href{https://github.com/fullscreen-triangle/lagrangian}{github.com/fullscreen-triangle/lagrangian}
}

\cvitem{borgia}{
  \textbf{Zero-Parameter Molecular Spectroscopy.}
  Derives electronic transitions, absorption, emission, and vibrational modes
  from bounded phase space geometry --- zero empirical parameters.
  Applicable to pre-screening linker chemistries and bioreceptor candidates:
  predict absorption cross-sections and refractive index contributions from
  molecular structure before running chip experiments.
  Validated against NIST: nine elements, six molecules; errors $<$1.63\,\%.
  Live: \href{https://honjo-masamune.vercel.app/tools}{honjo-masamune.vercel.app/tools}.
  \href{https://github.com/fullscreen-triangle/borgia}{github.com/fullscreen-triangle/borgia}
}

\cvitem{crown-prince + syndrome}{
  \textbf{Binding Affinity and Molecular Recognition.}
  Crown Prince: $K_d$ from bounded phase space geometry; Hill equation as
  limiting case; validated against 39 NIST compounds. Applicable to
  receptor--analyte $K_d$ estimation without empirical docking.
  Syndrome: spectral DFT molecular recognition ($\rho\geq0.72$ binding
  threshold; $O(cL\log L)$; same physics as label-free biosensor spectral
  complementarity). MHC-binding tool live at
  \href{https://syndrome-seven.vercel.app/tools/mhc-binding}{syndrome-seven.vercel.app}.
}

\cvitem{lavoisier}{
  \textbf{Spectral Data Analysis Pipeline.}
  Dual-pipeline analysis: numerical + CNN visual (PyTorch); 98.9\,\% NIST
  accuracy; Ray distributed; $>$100\,GB dataset support; memory-mapped I/O.
  Signal extraction methodology directly applicable to interferometric and
  resonance-shift biosensor readouts.
  \href{https://github.com/fullscreen-triangle/lavoisier}{github.com/fullscreen-triangle/lavoisier}
}

%----------------------------------------------------------------------------------------
%	TECHNICAL SKILLS
%----------------------------------------------------------------------------------------

\section{Technical Skills}

\cvitem{Biochemistry}{
  Protein purification and handling; antibody-based assay development (ELISA,
  lateral flow); bioconjugation and immobilisation strategies; MALDI sample
  surface preparation; digestion and sample prep protocols; reproducibility
  and QC in bioanalytical method development
}

\cvitem{Analytical Chemistry}{
  LC-MS/MS, MALDI-TOF, UPLC-MS/MS, TIMS-PASEF, QTOF;
  ion annotation (LIPID MAPS, LipidBlast, HMDB, NIST);
  spectral data analysis (numerical and CNN-based pipelines)
}

\cvitem{Computational / ML}{
  Python (primary), R, Rust; PyTorch; scikit-learn; Ray, Dask;
  Docker, FastAPI; first-principles optical and molecular physics modelling
}

%----------------------------------------------------------------------------------------
%	LANGUAGES
%----------------------------------------------------------------------------------------

\section{Languages}

\cvitemwithcomment{English}{Native / Fluent}{}
\cvitemwithcomment{German}{Conversational}{Living in Germany since 2011}

%----------------------------------------------------------------------------------------

\renewcommand{\listitemsymbol}{-~}

\end{document}
